% 期望风险的三个层次：贝叶斯风险、空间内最优风险与实际输出风险。
\begin{tikzpicture}[x=1mm,y=1mm,font=\kaishu\footnotesize]
  % 风险方向与三个风险水平
  \draw[-{Stealth[length=2mm]},line width=0.7pt,draw=BookMuted]
    (5,4) -- (5,46);
  \node[rotate=90,text=BookMuted,font=\kaishu\scriptsize]
    at (0,25) {期望风险升高};

  \node[anchor=east,font=\heiti\small,text=BookInk] at (17,10) {5\%};
  \draw[line width=0.9pt,draw=BookMuted] (20,10) -- (68,10);
  \node[anchor=south west,text=BookInk] at (20,11)
    {贝叶斯风险 $R_{\mathcal P}^\star$};
  \node[anchor=north west,text=BookMuted,font=\kaishu\scriptsize] at (20,9)
    {当前输入信息下的最低期望风险};

  \node[anchor=east,font=\heiti\small,text=BookInk] at (17,31) {12\%};
  \draw[line width=0.9pt,draw=BookAmber] (20,31) -- (68,31);
  \node[anchor=south west,text=BookInk] at (20,32)
    {类内最低期望风险 $\inf_{h\in\mathcal H}R_{\mathcal P}(h)$};

  \node[anchor=east,font=\heiti\small,text=BookInk] at (17,42) {15\%};
  \draw[line width=0.9pt,draw=BookBlue,densely dashed] (20,42) -- (68,42);
  \node[anchor=south west,text=BookInk] at (20,43)
    {输出假设的期望风险 $R_{\mathcal P}(\hat h)$};

  % 两段差距分别对应近似误差与估计误差；线型与文字共同编码语义。
  \draw[{Stealth[length=1.8mm]}-{Stealth[length=1.8mm]},
    line width=1pt,draw=BookAmber] (73,11) -- (73,30);
  \node[anchor=west,align=left,text width=30mm,text=BookAmber,font=\heiti\footnotesize]
    at (77,20.5) {近似误差：7个百分点};
  \node[anchor=west,align=left,text width=30mm,text=BookMuted,font=\kaishu\scriptsize]
    at (77,16.5) {假设空间没有包含更好的假设};

  \draw[{Stealth[length=1.8mm]}-{Stealth[length=1.8mm]},
    line width=1pt,draw=BookBlue,densely dashed] (73,32) -- (73,41);
  \node[anchor=west,align=left,text width=30mm,text=BookBlue,font=\heiti\footnotesize]
    at (77,38) {估计误差：3个百分点};
  \node[anchor=west,align=left,text width=30mm,text=BookMuted,font=\kaishu\scriptsize]
    at (77,34) {有限样本下没有选到空间内最佳假设};
\end{tikzpicture}
